%; whizzy document
% latex beamer presentation.
% platex, latex-beamer でコンパイルすることを想定。 

%     Tokyo Debian Meeting resources
%     Copyright (C) 2007 Junichi Uekawa

%     This program is free software; you can redistribute it and/or modify
%     it under the terms of the GNU General Public License as published by
%     the Free Software Foundation; either version 2 of the License, or
%     (at your option) any later version.

%     This program is distributed in the hope that it will be useful,
%     but WITHOUT ANY WARRANTY; without even the implied warranty of
%     MERCHANTABILITY or FITNESS FOR A PARTICULAR PURPOSE.  See the
%     GNU General Public License for more details.

%     You should have received a copy of the GNU General Public License
%     along with this program; if not, write to the Free Software
%     Foundation, Inc., 51 Franklin St, Fifth Floor, Boston, MA  02110-1301 USA

% 実行順番
% sudo  ~/bin/usb-macbook-ir.c &
% real presentation (shell-command (concat "DISPLAY=:0.1 xpdf -fullscreen " (replace-regexp-in-string "tex$" "pdf"(buffer-file-name)) "&"))
% DISPLAY=:0.1 xpdf -fullscreen 

\documentclass[cjk,dvipdfmx]{beamer}
\usetheme{Warsaw}
%  preview (shell-command (concat "xpdf " (replace-regexp-in-string "tex$" "pdf"(buffer-file-name)) "&"))
%  presentation (shell-command (concat "xpdf -fullscreen " (replace-regexp-in-string "tex$" "pdf"(buffer-file-name)) "&"))

%http://www.naney.org/diki/dk/hyperref.html
%日本語EUC系環境の時
\AtBeginDvi{\special{pdf:tounicode EUC-UCS2}}
%シフトJIS系環境の時
%\AtBeginDvi{\special{pdf:tounicode 90ms-RKSJ-UCS2}}

\title{東京エリア Debian 勉強会}
\subtitle{資料}
\author{上川 純一 dancer@debian.org\\IRC nick: dancerj}
\date{2007年2月17日}
\logo{\includegraphics[width=8cm]{image200607/openlogo-light.eps}}

% 三択問題用
\newcounter{santakucounter}
\newcommand{\santaku}[5]{%
\addtocounter{santakucounter}{1}
\frame{\frametitle{問題\arabic{santakucounter}. #1}
%問題\arabic{santakucounter}. #1
\begin{minipage}[t]{0.7\hsize}
 \begin{itemize}
 \item  A #2\\
 \item  B #3\\
 \item  C #4\\
 \end{itemize}
\end{minipage}
}
\frame{\frametitle{問題\arabic{santakucounter}. #1}
%問題\arabic{santakucounter}. #1
\begin{minipage}[t]{0.7\hsize}
\begin{itemize}
\item  A #2\\
\item  B #3\\
\item  C #4\\
\end{itemize}
\end{minipage}
\begin{minipage}[t]{0.2\hsize}
答えは:

\vspace{1cm}

{\huge \hspace{1cm}#5}
\end{minipage}}
}


\begin{document}
\frame{\titlepage{}}

\section{Intro}

\begin{frame}
 \frametitle{本日のagenda}
\begin{minipage}[t]{0.4\hsize}
  \begin{itemize}
  \item 注意事項
	\begin{itemize}
	 \item 飲食禁止
	 \item 政治/宗教/営利活動禁止
	\end{itemize}
  \item 最近事情
  \item 事前課題紹介
  \item quiz
 \end{itemize}
\end{minipage} 
\begin{minipage}[t]{0.4\hsize}
 \begin{itemize}
  \item quilt
  \item dbs 
  \item dpatch 
  \item OSC 2007 を企画する
 \end{itemize}
\end{minipage}
\end{frame}

\section{最近のDebian関連のミーティング報告}

\begin{frame}{最近のDebian関連ミーティング報告}
 \begin{itemize}
  \item Debian勉強会
 \end{itemize}
\end{frame}

\begin{frame}
 \frametitle{2007年1月のagenda}
\begin{minipage}[t]{0.4\hsize}
  \begin{itemize}
  \item 注意事項
	\begin{itemize}
	 \item 飲食禁止
	 \item 政治/宗教/営利活動禁止
	\end{itemize}
  \item 最近事情
  \item 事前課題紹介
  \item quiz
 \end{itemize}
\end{minipage} 
\begin{minipage}[t]{0.4\hsize}
 \begin{itemize}
  \item apt-torrent
  \item kvm
  \item 2007年のDebian勉強会を計画する
  \item quilt
  \item 今後のイベント
 \end{itemize}
\end{minipage}
\end{frame}

\section{事前課題の声}
\begin{frame}

 今回の事前課題は
「aptに足りなさそうな機能」と「パッケージングをしてみて感じたこと、または何故パッケージングをしないか」
 というタイトルで200-800文字程度の文章を書いてください。というものでした。
 その課題に対して下記の内容を提出いただきました。
\end{frame}


\subsection{やまねさん}
\begin{frame}{「aptに足りなさそうな機能」}

 \begin{itemize}
 \item 「aptのhttpではRedirectすらサポートしていない」らしいですよ？(\url{http://donrails.araki.net/archives/id/5580})
 \item apt というかフレームワークというか判りませんが、言語・環境ごとに規定してのインストール。メタパッケージでもできることはありますが、なんというか面倒な事も多いような。
 \end{itemize}
\end{frame}

\begin{frame}{「パッケージングをしてみて感じたこと、または何故パッケージングをしないか」}

 パッケージは作っているので、感じたことを。

 \begin{itemize}
 \item パズルみたいで面白い :-)
 \item プログラム書けなくても、まぁなんとかなるなる。
 \item distro specific なことをベタで書かれたものは扱いづらい。
 \end{itemize}
\end{frame}

\subsection{石原さん}
\begin{frame}
 {aptに足りなさそうな機能}

 apt のことをすべて把握しているわけじゃないのでなんともいえないです。
 使っているのは、
 \texttt{ apt-get install hogehoge}
 とか
 \texttt{ apt-get update;apt-get upgrade}
 くらいで、たまに
 \texttt{ apt-get remove hugahuga}
 とかです。

 こっちのへっぽこ設定で、時刻設定をハチャメチャなことにした状態で、
 apt-get をしてしまい、どうしようもなくなったことはありますが…
 \footnote{Windows Update の場合、ハチャメチャな時刻設定だと、たしかエラー
 になって出来なかったと思いますけど…。}

 あと、\texttt{apt-cache search} の結果（パッケージ説明）が
 日本語だといいなぁ、と思います。難しいのかなぁ？
\end{frame}
\begin{frame}{パッケージングをしてみて感じたこと、または何故パッケージングをしないか}

 パッケージングしたことないです。
 やり方がわかりません。
 そして、ホームディレクトリにパスを通して、
 \texttt{ ./configure ; make ; make install}
 で使えるし、ま、いっか、て感じなので調べていないですね…
\end{frame}

\subsection{澤田さん}

\begin{frame}{「パッケージングをしてみて感じたこと、または何故パッケージングをしないか」}

aptに足りなそうな機能とすると処理のロールバック機能だと思います。例えば、
\texttt{dist-upgrade} してみた→壊れてるパッケージがあって
\texttt{dist-upgrade}が完了できなかった→とりあえず壊れてるパッケージを
削除したくても中途半端な状態だと言われて削除できないorz→
\texttt{dist-upgrade}する前に戻れたら・・・といった感じです。ググってみ
るとそれっぽい試みはあるようですがaliothとかで引っ掛からないのだから形に
なっているものはないのかな？

\end{frame}

\begin{frame}{「パッケージングをしてみて感じたこと、または何故パッケージングをしないか」}

 パッケージングして感じたこととすると・・・、いい言葉が思いつかないので
 具体例を示します。うちのサーバには日本語パッチを当ててパッケージングし
 た\texttt{squirrelmail}が入っているのですが、DSAが出るたびに本家の最新とひとつ前
 をdiff→diffを見てローカルパッケージに手動パッチしてパッケージをビルド、
 しています。ここらへん、行番号は違っても回りから判断してパッチを当てて
 くれるツールがあるとうれしいなと思います。どっちかというとmore
 intelligent patchな話ですが本家配布のものをいじったパッケージを入れてるっ
 て場合に重要な気がします。調べたことないからひょっとしたらあるのかも。
 もしかして今回のネタそのもの？
\end{frame}

\subsection{キタハラさん}

\begin{frame}
 {aptに足りなさそうな機能}

 通常使用している範囲で、特に不足を感じる機能は
 ありません。
\end{frame}

\begin{frame}{パッケージングをしてみて感じたこと、または何故パッケージングをしないか}

 \begin{enumerate}
 \item  すみません、これまでパッケージングをした事
 がないので回答不能です。　

 \item  能力が無い＆これまで
 必要なモノは既にあり、必要がなかった。
 \end{enumerate}
\end{frame}

\subsection{えとーさん}

\begin{frame}
 何個か思い付きました。
 一個思い着くと連鎖でがんがんでるもんですな。

 \begin{enumerate}
 \item[1]  ケーパビリティの細分化とその明確化\\
 現状だとセキュアOSやLinuxケーパビリティを利用していてもパッケージ管理には
 フル権限が必要になる。
 パッケージ信頼性はMD5sumの提供や署名によりある程度は担保されるものの
 システム的な例外を作るしかない状況になっているので、なんとかしたい。
 運用で回避はできるのでまぁいいといえばいいのですが。。。コストもそれなりに掛るし
 美しくないので。
 \end{enumerate}
\end{frame}
\begin{frame}
\begin{enumerate}
 \item[2] ユーザ単位でのaptの利用\\
 現状ではrootユーザがシステムに対して行なうことが想定されているが
 それ以外にユーザがシステムに影響を与えずにインストールした場合もあるだろう。
 こいつをどうにかできないか。

 \item[3] パッケージの集中管理\\
 たとえサーバが二台しかなくても別々に見るの面倒なので一個のコンソールから
 管理したい。
 全ミラーじゃなくてインストールしているパッケージのみを選択的に
 ミラーしポリシーに従って自動的にアップデートを行なう。

 \end{enumerate}
\end{frame}
\begin{frame}
\begin{enumerate}

 \item[4] セキュリティアップデートのタグ付け\\
 サーバの場合はセキュリティアップデートとはいえ無造作に当てることはできない
 こともある。（できる場合は当然それでよいのだが）
 お仕事できっちりやる必要がある場合はアップデートがあればその差分ソースを
 精査し実験環境で動作確認を行なった上でアップデートを実行する必要があり、
 その場合はどうしようもないとしても、場合によっては管理コストから自動化したいし
 それが認められる状況（条件）もある。
 これをサポートするためにセキュリティアップデートにタグで情報の付加を行ない
 事前に選択した条件にマッチした場合のみアップデートを行なうという機能があれば
 便利ではなかろうか。
 また、バグの内容を解析して自動的に重みづけを行なうことや事前に重みづけを行なった
 機能との相関分析を行なうことによりアップデート条件の選択を簡易にできれば
 より堕落できるであろう。

 \end{enumerate}
\end{frame}


\subsection{あけどさん}

\begin{frame}
 {aptに足りなさそうな機能}

 aptを使っていてこんな機能が欲しいってのを挙げてみます。
 欲しい機能を持ったパッケージを探すのにapt-cache searchとかで検索することがあるのですが
 grep とかで工夫しても見つけるのに手間がかかるのを何とかしたいです。
 で、考えてみました。
 まずは、

 　Google とかの検索エンジンと連動する機能(仕組み?)  があればすごく便利に
 なるかと思います。今のapt系の検索ではライブラリとかアプリケーションなど
 のカテゴライズが今ひとつ整備できていない様に見えますので、それを補完する？
 機能が欲しいです。
\end{frame}

\begin{frame}
 もう一つ挙げてみますと、　パッケージをインストールせずにマニュアルページ
 だけ参照とか、パッケージ内の任意のファイルを参照(抽出)する機能ってのが
 aptで出来る様になれば良いかなと思います。

 以上の機能ってのはdpkgとか他のコマンドの組み合わせで出来るかもしれませんが
 まとまった物としてあれば良いと考えています。
\end{frame}

\subsection{前田 耕平さん}

\begin{frame}
 {「aptに足りなさそうな機能」}
\begin{enumerate}
 
 \item 足りなさそうな、というより実際足りない機能ですが、Vine
 Linuxのaptだと、apt-lineに書かずにローカルファイルも「apt-get install
 パッケージファイル名」でインストールできるのに、
 Debianではできないことですね。dpkgを使えば、と言われればそれまでなのですが、

 \item apt、というよりdpkgの方が妥当かと思いますが、rpmだと、-q
 --changelogでパッケージのChangeLogを表示できるのに、DebパッケージだとオプションでChangeLogを表示できないことでしょうか。（あるいは単に私が知らないだけ？）
\end{enumerate}
\end{frame}

\begin{frame}
 {「パッケージングをしてみて感じたこと、または何故パッケージングをしないか」}

 debianディレクトリ以下のファイルの編集が大変だなぁという感じがします。
 他には、普通にコンパイルするのはうまくいくのに、パッケージにするとうまくいかなかったりとか。
\end{frame}

\subsection{小室 文さん}

\begin{frame}
 {aptに足りなさそうな機能}

 \texttt{/etc/apt/source.list}を触る機会はあまりないですけど、最初(インス
 トール時) に\texttt{apt-spy}みたいなのが自動で動いてくれて
 \texttt{sources.list}に欲しいversionのミラーサイト書き込んでくれたらイン
 ストールの時に便利だなと思います。

 雑誌とか本についているインストールCDから入ると、設定していないのでそこで
 ミラーサイトなんだっけ？と検索したりそもそも記述方法が分からなかったり、
 とインストールorアップグレード時につまづく（かもしれない）要素を自分の過
 去を振り返って思いました。

 他の人はそうでもないのかもしれないけど・・・
\end{frame}


\subsection*{上川}
\begin{frame}[containsverbatim]{解答}

みなさまの不満は wajig でほとんど解決されているものですね。

\begin{texttt}
 Description: simplified Debian package management front end
 Wajig is a single commandline wrapper around apt, apt-cache, dpkg,
 /etc/init.d scripts and more, intended to be easy to use and providing
 extensive documentation for all of its functions.

 With a suitable sudo(1) configuration, most (if not all) package installation
 as well as creation tasks can be done from a user shell. Wajig is also
 suitable for general system administration.

 A Gnome GUI command 'gjig' is also included in the package.
\end{texttt}


\end{frame}

\section{DWNQuiz}
%% debianmeetingresume200702.texから以下コピー
\subsection{2007年01号}
\url{http://www.debian.org/News/weekly/2007/01/}
にある1月23日版です。

\santaku
{creative commons 2.5 は DFSG互換か？}
{ちがう}
{DFSG互換です}
{むしろGPLとまったく同じ}
{A}

\santaku
{Kenshi Mutoのアナウンスによると、Takeshi Yaegashiは何にDebianをインストールするのに成功した?}
{PLAYSTATION 3}
{Wii}
{XBox 360}
{A}

\santaku
{Florian Lohoffはどんな変化に気付いたか?}
{woodyがミラーから取り除かれた}
{sargeのアーカイブに侵入された形跡がある}
{etchへの開発者の興味が薄れている}
{A}

\santaku
{Joseph Smidtはリリースに関して何を提案したか?}
{unstableとtestingだけをサポートするようにして保守作業を楽にしよう}
{etchはリリースされないまま古くなりつつあるので適当にリリースしてlennyに全力を注ぎ込もう}
{そろそろDebianもXPやVistaといったよくわからない名前をリリースにつけるようにしよう}
{A}

\subsection{2007年02号}
\url{http://www.debian.org/News/weekly/2007/02/}
にある1月30日版です。

\santaku
{Roberto C. Sanchezが提案したのは何か}
{全員が自分の誕生日をDebianとして祝うべきだ}
{Debianのスクリーンショットのレポジトリを準備して、それぞれのパッケージ
に対応させる}
{Debianを全員使うべきだ}
{B}

\santaku
{Robert Millanが作成したwin32用プログラムは?}
{Windows Vistaを自動的にダウンロードし上書きインストールしてくれるインストーラ}
{ユーザの手間を省くために、Outlook Expressのアドレス帳に載っているアドレスをlists.debian.orgのメーリングリストすべてに自動登録してくれるプログラム}
{Debian Installerを自動的にダウンロード・起動してくれるランチャー}
{C}

\santaku
{1月31日に締め切られたのは?}
{第24回東京エリアDebian勉強会のblogでの報告}
{Debian etchベータ版のテスターの募集}
{Debian Conference 7のスポンサー希望者の参加登録}
{C}

\santaku
{Luis Matosがetchのカーネルに関して提案したのは?}
{kvmが利用できるように2.6.20を使おう}
{LinuxではなくkFreeBSDを使おう}
{ハードウェアサポートのためにポイントリリースでカーネルの更新を行おう}
{C}


\subsection{2007年03号}
\url{http://www.debian.org/News/weekly/2007/03/}
にある2月13日版です。

\santaku
{DebianのウェブサイトについてManoj Srivastavaが気付いたのは?}
{デザインがやや古びており見た目がイマイチ}
{投票ページのナビゲーションバーが長すぎる}
{朝起きたらぐるぐるマークが逆回転になっていた}
{B}

\santaku
{今年のDebian Project Leader選挙のアナウンスから4時間。最初にノミネートしてきたのは?}
{Junichi Uekawa}
{Gustavo Franco}
{Bill Gates}
{B}

\santaku
{Debian Conference 2008はどこの国で開催されることに決定したか?}
{イラク}
{アルゼンチン}
{パプアニューギニア}
{B}


\begin{frame}{quilt}
小林さん
\end{frame}

\begin{frame}{dbs}
岩松さん
\end{frame}

\section{dpatch}

\begin{frame}{agenda}
\begin{minipage}[t]{0.4\hsize}
 dpatch とはなにものか
\end{minipage}
\begin{minipage}[t]{0.4\hsize}
\begin{itemize}
 \item dpatch ツール紹介
 \item dpatch の使いかた紹介
 \item マニアックな使いかた
\end{itemize}
\end{minipage}
\end{frame}

\subsection{dpatchツール紹介}
\begin{frame}{dpatchツール紹介}

dpatch は Debianのソースパッチを管理するツールです。

Debianパッケージでは、ソースパッケージは以下の構成になってます。
\begin{description}
    \item[.orig.tar.gz] オリジナルのtarball
    \item[.diff.gz] Debianで作成した差分
    \item[.dsc] dpkg用制御ファイル
\end{description}

この中で、.diff.gzは一つの大きな差分ファイルとして管理されます。そのため、
どの部分がどういうパッチであるかは管理してはくれません。
そこで dpatch が活躍します。

通常の.diff.gzであると \texttt{debian/}ディレクトリ以下のDebianパッケー
ジング用の情報と それ以外のソフトウェア自体への修正が混合していますが、
それを整理してくれます。

\end{frame}

\begin{frame}{ソースのディレクトリ構成}

dpatch で管理しているソースのディレクトリ構成は下記のようになります。
 
\begin{description}
 \item[debian/patches/00list] 適用するパッチの一覧の記述したファイル
 \item[debian/patches/XXxxxx.dpatch] パッチファイル
 \item[debian/patched/] dpatch がパッチを適用したか管理するためのフラグファイル
\end{description}

\end{frame}

\begin{frame}
{使えるツール} 

一番よく利用するのは、dpatch-edit-patch でしょう。

\begin{description}
 \item[dpatch] メインのツールです。表立って直接利用することはありません。
	    debian/rules などから呼び出されます。
 \item[dpatch-list-patch] パッチの一覧を出力します。
 \item[dpatch-edit-patch] 編集のツールです。パッチを作成、もしくは編集する際に利用しま
	    す。
 \item[dpatch-convert-diffgz] .diff.gz から dpatch ファイルを生成するた
	    めのツールです。
 \item[dpatch-get-origtargz] dpatch-edit-patch が内部的に利用するツール。
\end{description}
\end{frame}

\subsection{dpatchの使いかた}
\begin{frame}{dpatchの使いかた}
\begin{enumerate}
 \item  \url{/usr/share/doc/dpatch/examples/rules/rules.new.dh.gz} を参照して、
 \texttt{ debian/rules} を書き換えます
 \item dpatch-edit-patch でパッチを作成します。
 \item debian/patches/00list を編集します
\end{enumerate}
\end{frame}

\subsection{マニアックな使いかた}
\begin{frame}{マニアックな使いかた}

\begin{itemize}
 \item  dbs のような使いかたもできる!

 \item  debian/ ディレクトリしか展開されていない場合でも。

 \item  デモ:\\
 cd demo/dpatch/
\end{itemize}

\end{frame}


\section{今後のイベント}

\begin{frame}
 \frametitle{今後の企画}
 \begin{itemize}
  \item 2007年3月 OSC
  \item 2007年4月 Debian 勉強会・インストール大会
  \item 2007年5月 Debian 勉強会
  \item 2007年6月 Debian Conference
  \item 2007年7月 Debian Conference 参加報告会
  \item 2007年7月 OSC関西？
 \end{itemize}
\end{frame}

\begin{frame}{OSC企画会議}
\begin{itemize}
 \item Debian勉強会と仮想化友の会で勉強会を開催
 \item 何ができるか?
 \item 何をする必要があるか?
\end{itemize}
\end{frame}

\begin{frame}{ターゲット}
\begin{minipage}[t]{0.4\hsize}
\begin{itemize}
 \item 仮想化使っている人
 \item 仮想化知っている人
 \item 仮想化知らない人
\end{itemize}
\end{minipage}
\begin{minipage}[t]{0.5\hsize}
\begin{itemize}
 \item Debianの開発者
 \item Debianのユーザ
       \begin{itemize}
	\item サーバ管理
	\item 開発者
	\item なぜかデスクトップユーザ
	\item その他
       \end{itemize}
 \item ユーザじゃない人
\end{itemize}
\end{minipage}
\end{frame}

\begin{frame}{前提}

想定している参加する人たちの比率予想（数字は適当）
\begin{tabular}{|c|c|c|c|}
\hline
 & DD & Debianユーザ& ユーザじゃない \\
\hline
仮想化使っている & 1\% & 25\% & 25\% \\
\hline
仮想化知っている & 0\% & 25\% & 25\% \\
\hline
仮想化知らない & 0\% & 0\% & 0\% \\
\hline
\end{tabular}
\end{frame}

\begin{frame}{テーマ}
Debianがホストだという前提。

仮想化ってどうやってつかうの?
\begin{itemize}[<+->]
 \item 開発・検証用
 \item セキュリティーのサンドボックス
 \item Windows を動かす/EDY チャージ、年賀状、winny
 \item 教育用・試験用
 \item テンプレート配布
 \item ?
\end{itemize}
\end{frame}

\begin{frame}{OSC出し物案}
\begin{minipage}[t]{0.4\hsize}
\begin{itemize}
 \item Introduction to Debian 勉強会
 \item やらせでSocial Contract 唱和
 \item やらせで事前課題をする
 \item 仮想化常識 Quiz (立ってください)
\end{itemize}
\end{minipage}
\begin{minipage}[t]{0.4\hsize}
\begin{itemize}
 \item Windows
 \item 仮想化技術紹介
 \item iSCSIでライブマイグレーションとか
 \item 最後の挨拶
%  \item パネルディスカッション
%        「仮想化時代のDebianディストリビューション(仮)」
%        \begin{itemize}
% 	\item 仮想化の各代表
% 	\item Debian の代表
%        \end{itemize}
\end{itemize}
\end{minipage}
\end{frame}

\begin{frame}
 \frametitle{宴会会場}
 \begin{itemize}
  \item 会場:荻窪 卯(うさぎ)
  \item 時間:21:00-
  \item 集合場所: この建物 1F のところ
  \item 注意事項
	\begin{itemize}
	 \item 終電の時間をちゃんとしらべていきましょう
	\end{itemize}
 \end{itemize}
\end{frame}

\end{document}
